\documentclass[12pt,a4paper]{article}
\usepackage[utf8]{inputenc}
\usepackage[vietnamese]{babel}
\usepackage[left=2cm,right=2cm,top=2cm,bottom=2cm]{geometry}
\usepackage{mathptmx}
\usepackage{enumitem}

\pagestyle{plain}
\linespread{1.15}

\begin{document}

% --- HEADER QUỐC HIỆU & CƠ QUAN BAN HÀNH ---
\noindent
\begin{minipage}[t]{0.45\textwidth}
    \begin{center}
        \small
        CÔNG AN THÀNH PHỐ HÀ NỘI\\
        \textbf{\underline{PHÒNG CẢNH SÁT HÌNH SỰ}}\\[0.1cm]
        Số: 15/BB-HC\\
        \textit{V/v hỏi cung bị can nhận tội}
    \end{center}
\end{minipage}
\hfill
\begin{minipage}[t]{0.52\textwidth}
    \begin{center}
        \small
        \textbf{CỘNG HÒA XÃ HỘI CHỦ NGHĨA VIỆT NAM}\\
        \textbf{\underline{Độc lập - Tự do - Hạnh phúc}}\\[0.1cm]
        \textit{Hà Nội, ngày 25 tháng 07 năm 2026}
    \end{center}
\end{minipage}

\vspace{0.6cm}

% --- TIÊU ĐỀ VĂN BẢN ---
\begin{center}
    \textbf{\large BIÊN BẢN HỎI CUNG BỊ CAN}\\
    \textit{(Trận đối chất bẻ gãy hoàn toàn lời khai gian dối của Trần Thị Hà)}
\end{center}

\vspace{0.3cm}

\noindent Vào hồi 14 giờ 30 phút ngày 25 tháng 07 năm 2026, tại Phòng Hỏi cung số 2, Trụ sở Phòng CSHS.\\
Chúng tôi gồm:
\begin{enumerate}[label=\arabic*., leftmargin=1.5cm]
    \item Điều tra viên cao cấp: Thượng tá Trần Quốc Dũng — Phó Trưởng phòng PC02.
    \item Cán bộ ghi biên bản: Trung úy Nguyễn Văn Hoàng — Cán bộ điều tra.
    \item Kiểm sát viên kiểm sát: Ông Nguyễn Văn Hòa — VKSND TP. Hà Nội.
\end{enumerate}

\noindent Tiến hành hỏi cung bị can: **TRẦN THỊ HÀ (Sinh ngày 22/09/1994)**.

\section*{NỘI DUNG ĐỐI CHẤT VÀ TỰ THÚ (HỎI VÀ ĐÁP)}

\noindent \textbf{Hỏi (Thượng tá Dũng):} Chị Hà, trong biên bản số 07d, chị khai rất trôi chảy: Tối qua sau khi ăn cơm xong, từ 20h00 đến 21h15 chị nằm trên giường tại phòng trọ xem phim bộ truyền hình VTV3 cho đỡ buồn, xem hết phim mới tắt tivi đi ngủ đúng không?\\
\textbf{Đáp (Trần Thị Hà):} *(Nhỏ nhẹ, vẻ mặt thành khẩn)* Vâng thưa cán bộ... tôi ở một mình buồn nên tối nào cũng bật phim bộ VTV3 xem quen rồi...

\noindent \textbf{Hỏi (Thượng tá Dũng):} Chị Hà à, chị ngụy tạo một bằng chứng ngoại phạm nghe rất bình dị, nhưng chị quên mất một điều cốt tử: **Tối qua là Thứ Sáu, ngày 24 tháng 7**. Khung giờ vàng tối cuối tuần trên VTV3 từ trước đến nay **luôn phát sóng chương trình gameshow thực tế, TUYỆT ĐỐI KHÔNG HỀ CHIẾU PHIM TRUYỀN HÌNH NÀO!** Bà Lụa ở nhà số 16 sát vách nhà Khang đã khai rất rõ trong Biên bản số 06: Tối qua bà ngồi xem gameshow thực tế trên VTV3! Vậy chị ngồi xem "bộ phim truyền hình" nào trên VTV3 vào tối Thứ Sáu?\\
\textbf{Đáp (Trần Thị Hà):} *(Sững người, đồng tử co lại, hai bàn tay đan chặt vào nhau)* Tôi... tôi xem tivi ngủ quên lúc nào không rõ... chắc là tôi nhớ nhầm...

\noindent \textbf{Hỏi (Thượng tá Dũng):} *(Bật máy tính bảng mở file ghi âm hộp thư thoại trích xuất từ điện thoại nạn nhân Khang)*\\
Mời chị nghe lại đoạn tin nhắn thoại chị gửi vào máy Khang lúc **20 giờ 32 phút tối qua**. Phía sau giọng nói dặn Khang ngủ sớm của chị lọt rất rõ **tiếng còi tàu hỏa diesel hú 2 hồi dài và tiếng chuông rào chắn đường sắt leng keng**. Phòng trọ của chị cách đường sắt hơn 1km, xung quanh nhà cao tầng san sát, không thể nào thu được tiếng còi tàu và chuông rào chắn rõ như vậy! Nơi duy nhất cách rào chắn 30 mét chính là ngay trước cổng nhà số 14 của Khang! Lúc 20h32, chị đang đứng ngay dưới gốc cây xoan rình mò trước cửa nhà Khang!\\
\textbf{Đáp (Trần Thị Hà):} *(Mặt cắt không còn giọt máu, toàn thân run rẩy)* Tôi... tôi...

\noindent \textbf{Hỏi (Thượng tá Dũng):} *(Đặt túi zip đựng lọn tóc dính máu thu giữ trong áo ngực của Hà và Báo cáo pháp y số 12 lên bàn)*\\
Lọn tóc mai dính máu này giám định ADN khẳng định 100\% là máu của Khang! Báo cáo pháp y số 12 kết luận: Khang tử vong lúc 21h00 do nhát đâm vào cổ, trong khi Tùng đã lên xe khách rời bến từ 20h15. Dấu vân tay ngón trỏ và ngón cái của chị in hằn trên mảnh thủy tinh hung khí. Người duy nhất ở cạnh Khang lúc 21h00 chính là chị!\\
\textbf{Đáp (Trần Thị Hà):} *(Im lặng vài giây, rồi khóe môi bỗng nhếch lên một nụ cười lạnh lẽo đến gai người)*\\
Các anh nghĩ vạch trần được tôi là tôi sợ sao? Tôi vào nhà thấy anh ấy nằm ngất, định lay anh ấy dậy... Nhưng đúng lúc đó điện thoại sáng lên tin nhắn con bé Yến Nhi rủ đi du lịch Đà Lạt, bảo giấu con mụ kế toán phiền phức là tôi đi... 

Tôi yêu anh ấy suốt 3 năm, chăm sóc từng sợi tóc, từng chiếc cúc áo của anh ấy... Vậy mà anh ấy dám phản bội tôi để đi với con đàn bà khác! Khang chết rồi... Anh ấy không thể đi Đà Lạt với ai được nữa. Mãi mãi anh ấy chỉ thuộc về một mình tôi...

\vspace{0.6cm}
\noindent Biên bản kết thúc hồi 16 giờ 00 phút cùng ngày.

\vspace{0.8cm}

% --- CHỮ KÝ NGHỊ ĐỊNH 30 ---
\noindent
\begin{minipage}[t]{0.45\textwidth}
    \small
    \textbf{\textit{Nơi nhận:}}\\
    - Viện KSMND TP. Hà Nội;\\
    - Thủ trưởng Cơ quan CSĐT;\\
    - Lưu: VT, CSHS.
\end{minipage}
\hfill
\begin{minipage}[t]{0.5\textwidth}
    \begin{center}
        \small
        \textbf{CÁN BỘ CHỦ TRÌ HỎI CUNG}\\
        \textit{(Ký, ghi rõ họ tên)}\\[1.5cm]
        \textbf{Thượng tá Trần Quốc Dũng}
    \end{center}
\end{minipage}

\end{document}
